\documentclass[12pt,letterpaper]{article}

% ----- Paquetes -----
\usepackage[utf8]{inputenc}
\usepackage[T1]{fontenc}
\usepackage[spanish]{babel}
\usepackage{amsmath,amssymb}
\usepackage{geometry}
\geometry{margin=2.8cm}
\usepackage{setspace}
\onehalfspacing
\usepackage{enumitem}
\usepackage[hidelinks]{hyperref}
\usepackage{booktabs}
\usepackage{array}
\usepackage{longtable}

% ----- Notación -----
\newcommand{\rv}{\mathbf{r}}
\newcommand{\uv}{\mathbf{u}}
\newcommand{\jv}{\mathbf{j}}
\newcommand{\Av}{\mathbf{A}}
\newcommand{\Bv}{\mathbf{B}}
\newcommand{\grad}{\nabla}
\newcommand{\gradp}{\nabla^{\perp}}
\newcommand{\EST}{\textsc{[Establecido]}}
\newcommand{\LIT}{\textsc{[Literatura]}}
\newcommand{\PRO}{\textsc{[Propuesto]}}

\title{\textbf{Análisis del problema inverso de la ecuación de Schrödinger\\
estacionaria con potencial vectorial:\\
formulación hidrodinámica y estudio de estados de Fock--Darwin}}
\author{Juan Sebastián [apellidos]\\
\small Director: John Hernán Díaz Forero\\
\small Proyecto Curricular de Física --- Universidad Distrital Francisco José de Caldas}
\date{\small Anteproyecto de Trabajo de Grado --- Modalidad Monografía de investigación teórica\\
(Acuerdo 004 de 2026, Art.~2 lit.~b; estructura del Art.~24)\\ Versión 2 --- \today}

\begin{document}
\maketitle

% =====================================================================
\section*{a) Tema o título}
\addcontentsline{toc}{section}{a) Tema o título}
El tema es el problema inverso de la ecuación de Schrödinger estacionaria para una partícula cargada en presencia de un
potencial vectorial: a partir de información sobre un estado estacionario, reconstruir los potenciales escalar y
vectorial que lo sostienen. Se delimita en contenido a la formulación hidrodinámica de Madelung y a las condiciones de
reconstrucción que de ella se derivan, incluida la libertad de gauge; en espacio, a sistemas bidimensionales con campo
magnético perpendicular; y en tiempo, al plazo del Art.~9 del Acuerdo 004 de 2026. Los estados de Fock--Darwin sirven
como caso de estudio de solución exacta conocida.

% =====================================================================
\section*{b) Introducción}
\addcontentsline{toc}{section}{b) Introducción}
La mecánica cuántica suele estudiarse en sentido directo: dado un potencial, se calculan los estados estacionarios. El
problema inverso pregunta lo contrario: qué potencial produce un estado o un espectro dado. Esa pregunta tiene una
tradición larga, que va de los problemas espectrales en una dimensión a los problemas de dispersión y de frontera con
campos magnéticos. Recientemente, Bhalgamiya y Novotny (2024) propusieron una versión constructiva que obtiene el
potencial directamente a partir de la densidad de probabilidad de un estado fundamental real.

Este trabajo estudia qué ocurre con esa construcción cuando el estado es complejo y la partícula está sometida a un campo
magnético. \textbf{Qué:} un análisis de las condiciones bajo las cuales se pueden reconstruir los potenciales escalar y
vectorial. \textbf{Cómo:} mediante la formulación hidrodinámica de Madelung, una revisión documental sistemática y un
estudio teórico-computacional de algunos estados de Fock--Darwin. \textbf{Por qué:} al introducir el potencial vectorial
aparecen la libertad de gauge y un problema de unicidad que no existen en el caso real. \textbf{Cuándo y dónde:} en el
plazo reglamentario, con herramientas de libre acceso y sin trabajo de campo. \textbf{Para qué:} para establecer qué
información es necesaria para que la reconstrucción sea posible y cuáles son sus limitaciones.

% =====================================================================
\section*{c) Situación problémica y pregunta de investigación}
\addcontentsline{toc}{section}{c) Situación problémica y pregunta}

\paragraph{Cómo surge el problema.} La revisión de antecedentes (sección e) muestra tres líneas de trabajo que
apenas se tocan:
\begin{enumerate}[leftmargin=1.8em]
  \item Los problemas inversos espectrales, de dispersión y de frontera reconstruyen potenciales a partir de datos
        \emph{externos}: espectros, matrices de dispersión o mediciones en la frontera. Con campo magnético, sus resultados
        de unicidad recuperan el campo $\Bv$ y el potencial escalar, pero el potencial vectorial solo salvo un gauge,
        porque los datos son invariantes de gauge.
  \item Las construcciones a partir del estado fundamental (mecánica cuántica supersimétrica; Bhalgamiya y Novotny)
        usan datos \emph{internos}, la función de onda o la densidad, pero suponen una función de onda real y sin campo
        magnético.
  \item La formulación hidrodinámica de Madelung y Bohm proporciona las ecuaciones para un estado complejo con potencial
        vectorial, pero se ha usado sobre todo para el problema directo.
\end{enumerate}
En la literatura revisada hasta ahora no se encontró un estudio de la intersección: la reconstrucción a partir de datos
internos de un estado complejo con potencial vectorial. En ese caso no está establecido qué datos internos determinan los
potenciales ni cómo interviene la libertad de gauge. Ese es el vacío que aborda el trabajo.

\paragraph{Delimitación física.}
\begin{center}
\begin{tabular}{@{}p{3.6cm} p{11cm}@{}}
\toprule
\textbf{Sistema físico} & Una partícula cargada (carga $q$, masa $m$), no relativista y sin espín, confinada en un plano
y sometida a un campo magnético perpendicular $\Bv=B_z(x,y)\,\hat z$, en un estado estacionario ligado y normalizable. \\
\textbf{Modelo} & Ecuación de Schrödinger estacionaria con acoplamiento mínimo,
$\frac{1}{2m}(-i\hbar\grad-q\Av)^2\psi+V\psi=E\psi$, con $\psi$ de clase $C^2$ fuera de un conjunto nodal de medida nula
y decaimiento en el infinito (o condiciones de Dirichlet en un dominio acotado). \\
\textbf{Objeto de estudio} & La reconstrucción del potencial escalar $V$ y del potencial vectorial $\Av$ (o del campo
$B_z$) y las condiciones bajo las cuales es posible y única. \\
\textbf{Magnitudes} & Densidad $P=|\psi|^2$; fase $S$; corriente $\jv=P(\grad S-q\Av)/m$; campo $B_z$; energía $E$;
condiciones de frontera. \\
\bottomrule
\end{tabular}
\end{center}

\paragraph{El problema inverso: datos, incógnitas e información adicional.} Con $\psi=\sqrt P\,e^{iS/\hbar}$ y
$\uv=\grad S-q\Av$, la ecuación de Schrödinger se separa en la continuidad estacionaria y la ecuación maestra extendida:
\begin{align}
  \grad\cdot(P\uv)&=0, \label{eq:cont}\\
  V-E&=-\frac{1}{2m}|\uv|^2+\frac{\hbar^2}{4m}\!\left[\frac{\grad^2P}{P}-\frac12\frac{|\grad P|^2}{P^2}\right].
  \label{eq:maestra}
\end{align}
Según qué se tome como dato, el problema cambia de naturaleza:
\begin{center}
\begin{tabular}{@{}p{2.6cm} p{3.2cm} p{8.6cm}@{}}
\toprule
\textbf{Datos} & \textbf{Incógnitas} & \textbf{Qué se determina} (resultados preliminares, por demostrar con rigor) \\
\midrule
$(P,\jv)$ & $V$, $\Av$ & $\uv=m\jv/P$; $B_z=-\frac{1}{q}(\grad\times\uv)_z$ donde $S$ es regular, y $V-E$ por~\eqref{eq:maestra}.
$\Av$ solo salvo gauge, como en los problemas de frontera. La continuidad es una condición de compatibilidad sobre los datos. \\
$(P,S)$ en un gauge & $V$, $\Av$ & La continuidad es una sola ecuación para $\Av$ y \emph{no} determina $B_z$: queda libre una
función arbitraria de $P$ (sección d.3). \\
$(P,S,B_z)$ & $V$, $\Av$ & $\Av$ queda fijado salvo $\grad\chi$, y $\chi$ resuelve una ecuación elíptica; luego $V$ por~\eqref{eq:maestra}. \\
$(P,S,V)$ & $\Av$ & \eqref{eq:maestra} fija $|\uv|$ y la continuidad restringe su dirección; problema no lineal por estudiar. \\
\bottomrule
\end{tabular}
\end{center}
En todos los casos $E$ entra como una constante aditiva en $V$. En los nodos de $\psi$ (núcleos de vórtice si $\ell\neq0$,
nodos radiales si $n>0$) el término de Bohm y la fase son singulares, y el rotacional de $\uv$ contiene contribuciones
concentradas en los vórtices, proporcionales a la circulación cuantizada $2\pi\hbar\ell$.

\medskip
\noindent\textbf{Pregunta de investigación.} ¿Bajo qué condiciones es posible reconstruir los potenciales asociados a una
partícula cargada descrita por la ecuación estacionaria de Schrödinger en dos dimensiones, a partir de la formulación
hidrodinámica de Madelung, y cuáles son las limitaciones de dicha reconstrucción en estados de Fock--Darwin?

% =====================================================================
\section*{d) Marco teórico}
\addcontentsline{toc}{section}{d) Marco teórico}
Para distinguir el estatus de cada elemento se usan tres etiquetas: \EST{} (fundamento teórico aceptado), \LIT{}
(resultado reportado en la literatura) y \PRO{} (desarrollo que el trabajo propone realizar o demostrar).

\subsection*{d.1 Acoplamiento mínimo e invariancia de gauge \EST}
La ecuación de la sección c es invariante bajo $\Av\to\Av+\grad\chi$, $\psi\to e^{iq\chi/\hbar}\psi$, es decir,
$S\to S+q\chi$. Son invariantes $P$, $\uv$, $\jv$, $V-E$ y $\Bv=\grad\times\Av$. En dominios no simplemente conexos (por
ejemplo, el plano sin los nodos de $\psi$) los potenciales pueden tener efectos que no se reducen al campo local
(Aharonov y Bohm, 1959).

\subsection*{d.2 Formulación hidrodinámica \EST}
Bajo las condiciones del modelo, la parte imaginaria de la ecuación da~\eqref{eq:cont} y la real da
$V-E=\frac{\hbar^2}{2m}\frac{\grad^2\sqrt P}{\sqrt P}-\frac{1}{2m}|\uv|^2$; con la identidad
$\grad^2\sqrt P/\sqrt P=\frac12\grad^2P/P-\frac14|\grad P|^2/P^2$ se obtiene~\eqref{eq:maestra} (Madelung, 1927; Bohm,
1952; Takabayasi, 1952). La fase es multivaluada alrededor de los nodos: la circulación de $\grad S$ alrededor de un nodo
simple es $2\pi\hbar\ell$ (Dirac, 1931).

\subsection*{d.3 Reconstrucción a partir del estado fundamental real \LIT}
Para $\psi$ real y positiva, $\Av=\mathbf 0$ y $S$ constante, \eqref{eq:cont} se cumple idénticamente y~\eqref{eq:maestra}
con $P=e^{-f}$ da $V-E=-\frac{\hbar^2}{4m}\grad^2f+\frac{\hbar^2}{8m}|\grad f|^2$ (Bhalgamiya y Novotny, 2024). En una
dimensión equivale a la construcción del superpotencial de la mecánica cuántica supersimétrica,
$W\propto-\psi_0'/\psi_0$ (Cooper, Khare y Sukhatme, 1995).

\subsection*{d.4 Unicidad salvo gauge en problemas inversos magnéticos \LIT}
Con datos de frontera (mapa de Dirichlet a Neumann) o de dispersión a energía fija, el operador de Schrödinger magnético
determina $\Bv$ y $V$, pero no $\Av$, por la invariancia de gauge de los datos (Sun, 1993; Nakamura, Sun y Uhlmann, 1995;
Eskin y Ralston, 1995). Los resultados citados se formulan, en general, en dimensión mayor o igual a tres; el caso
bidimensional requiere técnicas distintas (por precisar en la revisión).

\subsection*{d.5 Estados de Fock--Darwin \LIT}
Para $V=\frac12m\omega_0^2r^2$ y $\Bv=B\hat z$ en el gauge simétrico $\Av=\frac{Br}{2}\hat\theta$, con $\omega_c=qB/m$,
$\Omega=\sqrt{\omega_0^2+\omega_c^2/4}$ y $l_\Omega=\sqrt{\hbar/(m\Omega)}$:
\begin{equation}
  \psi_{n,\ell}\propto\Big(\tfrac{r}{l_\Omega}\Big)^{|\ell|}L_n^{|\ell|}\!\Big(\tfrac{r^2}{l_\Omega^2}\Big)e^{-r^2/2l_\Omega^2}e^{i\ell\theta},
  \qquad E_{n,\ell}=\hbar\Omega(2n+|\ell|+1)-\tfrac12\hbar\omega_c\ell
\end{equation}
(Fock, 1928; Darwin, 1931). Son el modelo estándar de un electrón en un punto cuántico parabólico y conectan, al crecer
$\omega_c$, con los niveles de Landau (Reimann y Manninen, 2002).

\subsection*{d.6 Desarrollos que propone el trabajo \PRO}
\begin{enumerate}[leftmargin=1.8em]
  \item \textbf{Reconstrucción con datos invariantes.} Demostrar que $(P,\jv)$ determinan $B_z$ y $V-E$ fuera de los nodos,
        y precisar qué ocurre en los nodos.
  \item \textbf{Degeneración con datos de fase.} En el gauge de Coulomb, $\Av=\gradp\Phi$ con $\gradp\Phi=(\partial_y\Phi,-\partial_x\Phi)$
        y $B_z=-\grad^2\Phi$, la continuidad se reduce a $q\{P,\Phi\}=\grad\cdot(P\grad S)$, con
        $\{P,\Phi\}=\partial_xP\,\partial_y\Phi-\partial_yP\,\partial_x\Phi$. Si $\Phi$ es solución, también lo es
        $\Phi+F(P)$ para toda $F$, mientras $B_z$ cambia en $-\grad^2F(P)$: los datos $(P,S)$ no determinan el campo.
        La degeneración en simetría axial es el caso particular $F(P(r))$. Esto se verificó numéricamente con una densidad
        sin simetría; la demostración general y la condición de solubilidad sobre curvas de nivel cerradas son parte del
        trabajo.
  \item \textbf{Comportamiento en los nodos.} Analizar la cancelación entre el término centrífugo de $|\uv|^2$ y la parte
        singular del término de Bohm cerca de los núcleos de vórtice, y el salto de fase $\pi\hbar$ en los nodos radiales.
\end{enumerate}

% =====================================================================
\section*{e) Antecedentes: trayectoria del problema}
\addcontentsline{toc}{section}{e) Antecedentes}

\paragraph{1. Problemas inversos espectrales en una dimensión (1929--1955).} Ambarzumian (1929) mostró que un espectro
particular determina el potencial en un caso especial; Borg (1946) estableció que en general hacen falta dos espectros;
Gel'fand y Levitan (1951) dieron un procedimiento de reconstrucción a partir de la función espectral, y Marchenko lo
desarrolló en paralelo. \emph{Aporte:} la reconstrucción es posible y única con datos espectrales suficientes.
\emph{Límite:} una dimensión, sin campo magnético.

\paragraph{2. Problemas inversos de dispersión (1950--1990).} La teoría se extendió a datos de dispersión; Chadan y
Sabatier (1977, 2.\textsuperscript{a} ed. 1989) la sistematizaron. \emph{Aporte:} métodos para varias dimensiones y
energías. \emph{Límite:} aparecen ambigüedades, potenciales distintos con los mismos datos, y los datos son externos al
sistema.

\paragraph{3. Problemas inversos con campo magnético (1993--1995 en adelante).} Sun (1993) observó que el mapa de
Dirichlet a Neumann es invariante de gauge y, por tanto, solo puede determinar $\Bv$ y $V$; Nakamura, Sun y Uhlmann (1995)
probaron la identificabilidad global y Eskin y Ralston (1995) el resultado análogo para dispersión a energía fija.
\emph{Aporte:} la unicidad salvo gauge es el resultado correcto, no una limitación. \emph{Límite:} datos de frontera o
de dispersión, y dimensión $\ge3$ en los resultados citados.

\paragraph{4. Construcción a partir del estado fundamental (1980--2024).} En una dimensión, el potencial se obtiene
directamente de un estado fundamental sin nodos mediante el superpotencial (Cooper, Khare y Sukhatme, 1995). Bhalgamiya
y Novotny (2024) generalizan la idea a 1D--3D con la densidad como dato, con aplicaciones al diseño de potenciales en
condensados de Bose--Einstein (trabajo previo: arXiv:2110.12082; tesis doctoral, Mississippi State University, 2023).
\emph{Aporte:} datos internos y fórmulas explícitas. \emph{Límite:} función de onda real; no incluye potencial vectorial
ni estados con corriente.

\paragraph{5. Formulación hidrodinámica (1927 en adelante).} Madelung (1927) escribió la ecuación de Schrödinger como un
fluido; Bohm (1952) y Takabayasi (1952) la desarrollaron, incluido el potencial cuántico; Dirac (1931) estableció la
cuantización de la circulación en los nodos, y la formulación se ha usado ampliamente en dinámica con trayectorias
(Holland, 1993; Wyatt, 2005). \emph{Aporte:} las ecuaciones~\eqref{eq:cont}--\eqref{eq:maestra} para estados complejos
con potencial vectorial. \emph{Límite:} uso casi exclusivo en el problema directo.

\paragraph{6. Estados de Fock--Darwin (1928 en adelante).} Fock (1928) y Darwin (1931) resolvieron el oscilador
bidimensional en campo magnético, que se volvió el modelo de referencia de los puntos cuánticos (Reimann y Manninen,
2002). \emph{Aporte:} un caso de solución exacta con vórtices, nodos radiales y un parámetro que interpola hacia el
régimen de Landau. \emph{Uso aquí:} banco de pruebas.

\paragraph{Comparación.}
\begin{center}
\begin{tabular}{@{}p{3.1cm} p{2.6cm} p{1.6cm} p{3.6cm} p{3.2cm}@{}}
\toprule
\textbf{Línea} & \textbf{Datos} & \textbf{$\Av$} & \textbf{Qué se determina} & \textbf{Abierto} \\
\midrule
1. Espectral 1D & espectros & no & $V$ único (dos espectros) & más dimensiones \\
2. Dispersión & matriz $S$ & no & $V$, con ambigüedades & unicidad general \\
3. Magnético & frontera, $S$ & sí & $\Bv$, $V$; $\Av$ salvo gauge & datos internos; 2D \\
4. Estado fundamental & $\psi$ o $P$ & no & $V-E$ explícito & estados complejos \\
5. Hidrodinámica & (directo) & sí & ecuaciones & uso inverso \\
6. Fock--Darwin & (exacto) & sí & solución completa & banco de pruebas \\
\bottomrule
\end{tabular}
\end{center}
La comparación ubica el vacío: las líneas 3 y 4 resuelven, por separado, el papel del gauge con datos externos y la
reconstrucción con datos internos, pero no el caso en que ambos aspectos aparecen a la vez. La línea 5 da las ecuaciones
para estudiarlo y la línea 6, el caso de prueba.

% =====================================================================
\section*{f) Justificación}
\addcontentsline{toc}{section}{f) Justificación}
\emph{Pertinencia física y matemática:} el trabajo precisa qué información sobre un estado cuántico con campo magnético
permite reconstruir los potenciales, y conecta un método constructivo reciente con los resultados clásicos de unicidad
salvo gauge. \emph{Relevancia formativa:} el problema integra temas centrales del plan de estudios (potencial vectorial,
invariancia de gauge, vórtices cuánticos, problemas mal condicionados y métodos numéricos) y produce material útil para
la docencia. \emph{Relevancia aplicada:} el diseño inverso de confinamientos es de interés en puntos cuánticos y en
potenciales ópticos para átomos fríos. \emph{Dimensión reflexiva} (complementaria): el caso permite discutir, a partir de
Barad (2007) y Latour (1987), cómo la elección de los datos, del gauge y de la regularización participa en lo que se
considera ``el potencial reconstruido''; esta discusión se incluirá como cierre del documento y no constituye un objetivo
independiente.

% =====================================================================
\section*{g) Objetivos}
\addcontentsline{toc}{section}{g) Objetivos}
\paragraph{h) Objetivo general.} Analizar las condiciones de reconstrucción de los potenciales escalar y vectorial en la
ecuación estacionaria de Schrödinger, mediante la formulación hidrodinámica de Madelung y el estudio teórico-computacional
de estados de Fock--Darwin.

\paragraph{i) Objetivos específicos.}
\begin{enumerate}[label=OE\arabic*.,leftmargin=3em]
  \item Sistematizar los fundamentos teóricos y los antecedentes bibliográficos del problema inverso de la ecuación de
        Schrödinger, particularmente en presencia de potenciales vectoriales.
  \item Derivar y analizar las ecuaciones de la formulación hidrodinámica de Madelung, identificando las condiciones de
        reconstrucción de los potenciales y el papel de la libertad de gauge.
  \item Estudiar analítica y numéricamente algunos casos representativos de estados de Fock--Darwin, con el propósito de
        examinar la reconstrucción de los potenciales y sus limitaciones.
  \item Contrastar los resultados obtenidos con los reportados en la literatura, identificando las condiciones bajo las
        cuales resulta aplicable el procedimiento de reconstrucción.
\end{enumerate}

% =====================================================================
\section*{j) Metodología}
\addcontentsline{toc}{section}{j) Metodología}
Investigación teórica con tres componentes: documental, analítico y computacional acotado.

\paragraph{Fase 1. Investigación documental (OE1).}
\begin{itemize}[leftmargin=1.6em]
  \item \emph{Búsqueda.} Bases: Web of Science, Scopus, arXiv, revistas de APS y AIP (incluida \emph{American Journal of
        Physics}) y Google Scholar. Cadenas de ejemplo: ``inverse problem'' AND Schrödinger AND (``magnetic potential'' OR
        ``vector potential''); ``inverse spectral problem''; ``Madelung'' OR ``quantum hydrodynamics''; ``Fock--Darwin''
        AND (``quantum dot'' OR ``inverse''); ``gauge'' AND ``uniqueness'' AND ``magnetic Schrödinger''.
  \item \emph{Bola de nieve.} Hacia atrás, desde las referencias de Bhalgamiya y Novotny (2024), Nakamura, Sun y Uhlmann
        (1995) y Cooper, Khare y Sukhatme (1995); hacia adelante, desde los trabajos que los citan.
  \item \emph{Selección.} Inclusión: trabajos que formulen un problema inverso para la ecuación de Schrödinger, traten la
        unicidad o el gauge, o desarrollen la formulación hidrodinámica o los estados de Fock--Darwin. Exclusión:
        aplicaciones a sistemas no cuánticos y trabajos sin resultados verificables. Se conservan las obras fundacionales
        sin restricción de fecha.
  \item \emph{Clasificación y sistematización.} Ficha por fuente en un gestor bibliográfico y matriz con las categorías:
        tipo de problema inverso, datos, incógnitas, presencia y tratamiento del potencial vectorial, dimensión, resultado
        de unicidad, método y limitaciones.
  \item \emph{Comparación.} Tablas comparativas como la de la sección e, a partir de las cuales se formula y se actualiza
        el vacío de conocimiento.
\end{itemize}

\paragraph{Fase 2. Desarrollo analítico (OE2).} Derivación de~\eqref{eq:cont}--\eqref{eq:maestra} con hipótesis
explícitas; análisis de los casos de datos de la tabla de la sección c; demostración de la degeneración con datos de fase
y de la reconstrucción con datos invariantes; análisis de las singularidades en los nodos.

\paragraph{Fase 3. Estudio de casos representativos (OE3).} Tres estados de Fock--Darwin con $\omega_c/\omega_0=2$:
$(n,\ell)=(0,0)$ como control sin vórtice, $(0,1)$ con vórtice y $(1,0)$ con nodo radial. Para cada uno: verificación
analítica de la reconstrucción, cálculo numérico en malla por diferencias finitas y estudio de convergencia; para el
estado $(0,1)$, además, sensibilidad al ruido en $P$. Se usa Python con bibliotecas de libre acceso, y cada prueba se
registra con su criterio de éxito antes de ejecutarla.

\paragraph{Fase 4. Contraste y redacción (OE4).} Comparación de los resultados con los de la literatura (líneas 3 y 4 de
la sección e) y redacción del informe final según el Art.~25 (RAE, resultados, conclusiones, recomendaciones y anexos).

\paragraph{Cronograma (dos periodos académicos).}
\begin{center}
\begin{tabular}{@{}p{3.1cm} p{6.6cm} p{1.4cm} p{3.2cm}@{}}
\toprule
\textbf{Periodo} & \textbf{Actividad} & \textbf{Obj.} & \textbf{Entregable} \\
\midrule
P1, sem.~1--6 & Investigación documental: búsqueda, selección y matriz & OE1 & Estado del arte \\
P1, sem.~5--14 & Desarrollo analítico & OE2 & Capítulo teórico \\
P1, sem.~13--16 & Casos $(0,0)$ y $(0,1)$ & OE3 & Primeros resultados \\
P2, sem.~1--6 & Caso $(1,0)$; convergencia y ruido & OE3 & Capítulo de resultados \\
P2, sem.~7--11 & Contraste con la literatura & OE4 & Capítulo de discusión \\
P2, sem.~12--16 & Redacción final, RAE y sustentación & todos & Documento y sustentación \\
\bottomrule
\end{tabular}
\end{center}
Si el plan de estudios contempla un solo espacio de Trabajo de Grado (Art.~9), se comprime a un periodo reduciendo la
Fase 3 a los estados $(0,0)$ y $(0,1)$.

% =====================================================================
\section*{k) Impacto esperado}
\addcontentsline{toc}{section}{k) Impacto esperado}
Una caracterización de qué datos permiten reconstruir los potenciales escalar y vectorial de un estado estacionario
bidimensional, contrastada con los resultados clásicos de unicidad salvo gauge; un estudio cuidadoso de tres estados de
Fock--Darwin que ilustra la reconstrucción y sus limitaciones (nodos, ruido, mallas); y un documento con valor
pedagógico. Como producto derivado posible, un artículo para una revista de física docente o computacional.

% =====================================================================
\section*{l) Bibliografía}
\addcontentsline{toc}{section}{l) Bibliografía}
\small Las entradas marcadas con $\dagger$ se citan de memoria y deben verificarse contra la fuente antes de la entrega.
\begin{enumerate}[leftmargin=2em,label={[\arabic*]}]
\item[] \textbf{Problemas inversos}
  \item V. A. Ambarzumian, ``Über eine Frage der Eigenwerttheorie'', \emph{Z. Phys.} \textbf{53}, 690 (1929).$\dagger$
  \item G. Borg, ``Eine Umkehrung der Sturm--Liouvilleschen Eigenwertaufgabe'', \emph{Acta Math.} \textbf{78}, 1 (1946).$\dagger$
  \item I. M. Gel'fand y B. M. Levitan, ``On the determination of a differential equation from its spectral function'',
        \emph{Izv. Akad. Nauk SSSR Ser. Mat.} \textbf{15}, 309--360 (1951).
  \item K. Chadan y P. C. Sabatier, \emph{Inverse Problems in Quantum Scattering Theory}, 2.\textsuperscript{a} ed.,
        Springer (1989).
  \item Z. Sun, ``An inverse boundary value problem for Schrödinger operators with vector potentials'', \emph{Trans. Amer.
        Math. Soc.} (1993).$\dagger$ (volumen y páginas)
  \item G. Nakamura, Z. Sun y G. Uhlmann, ``Global identifiability for an inverse problem for the Schrödinger equation in a
        magnetic field'', \emph{Math. Ann.} \textbf{303}, 377--388 (1995).
  \item G. Eskin y J. Ralston, ``Inverse scattering problem for the Schrödinger equation with magnetic potential at a fixed
        energy'', \emph{Commun. Math. Phys.} \textbf{173}, 199--224 (1995).
  \item F. Cooper, A. Khare y U. Sukhatme, ``Supersymmetry and quantum mechanics'', \emph{Phys. Rep.} \textbf{251},
        267--385 (1995).
  \item B. Bhalgamiya y M. A. Novotny, ``Exact solutions for the inverse problem of the time-independent Schrödinger
        equation'', \emph{Am. J. Phys.} \textbf{92}(12), 975--979 (2024).
  \item ``Example exact solutions of the time-independent Gross--Pitaevskii and Schrödinger equations'',
        arXiv:2110.12082.$\dagger$ (autores)
  \item B. Bhalgamiya, tesis doctoral, Mississippi State University (2023).$\dagger$ (título completo)
\item[] \textbf{Formulación hidrodinámica y gauge}
  \item E. Madelung, ``Quantentheorie in hydrodynamischer Form'', \emph{Z. Phys.} \textbf{40}, 322 (1927).$\dagger$
  \item P. A. M. Dirac, ``Quantised singularities in the electromagnetic field'', \emph{Proc. R. Soc. A} \textbf{133}, 60 (1931).$\dagger$
  \item D. Bohm, ``A suggested interpretation of the quantum theory in terms of hidden variables'', \emph{Phys. Rev.}
        \textbf{85}, 166 (1952).$\dagger$
  \item T. Takabayasi, ``On the formulation of quantum mechanics associated with classical pictures'', \emph{Prog. Theor.
        Phys.} \textbf{8}, 143 (1952).
  \item Y. Aharonov y D. Bohm, ``Significance of electromagnetic potentials in the quantum theory'', \emph{Phys. Rev.}
        \textbf{115}, 485 (1959).$\dagger$
  \item P. R. Holland, \emph{The Quantum Theory of Motion}, Cambridge University Press (1993).$\dagger$
  \item R. E. Wyatt, \emph{Quantum Dynamics with Trajectories}, Springer (2005).$\dagger$
\item[] \textbf{Estados de Fock--Darwin}
  \item V. Fock, ``Bemerkung zur Quantelung des harmonischen Oszillators im Magnetfeld'', \emph{Z. Phys.} \textbf{47}, 446 (1928).$\dagger$
  \item C. G. Darwin, ``The diamagnetism of the free electron'', \emph{Proc. Camb. Phil. Soc.} \textbf{27}, 86 (1931).$\dagger$
  \item S. M. Reimann y M. Manninen, ``Electronic structure of quantum dots'', \emph{Rev. Mod. Phys.} \textbf{74}, 1283 (2002).
\item[] \textbf{Dimensión reflexiva}
  \item K. Barad, \emph{Meeting the Universe Halfway}, Duke University Press (2007).
  \item B. Latour, \emph{Science in Action}, Harvard University Press (1987).
\item[] \textbf{Normativa}
  \item Universidad Distrital Francisco José de Caldas, Consejo Académico, Acuerdo N.\textsuperscript{o} 004 (18 de agosto de 2026).
\end{enumerate}

\vfill
\noindent\rule{\linewidth}{0.4pt}\\
\footnotesize Las comprobaciones numéricas preliminares están en el repositorio del proyecto
(\texttt{NucleoTG/teoria/verificar\_gauge\_fock\_darwin.py} y \texttt{NucleoTG/analisis/b1\_validacion.py}).

\end{document}
